% !TEX root = /home/eb/.vscode/extensions/qftrules.programme-colle/tmp/Exercice.tex


\begin{exo}[PC][3][TD]{La molécule d'ammoniac}
  \Notions{\item Puits de potentiel fini \task Double puits \task État stationnaire \task Effet tunnel}
  La molécule d'ammoniac $\mathrm{NH}_3$ est constituée d'un atome d'azote et de trois atomes d'hydrogène. L'atome d'azote peut se trouver au-dessus ou en-dessous du plan formé par les trois atomes d'hydrogène, ce qui correspond à deux configurations de la molécule. En modélisant la position de l'atome d'azote par une coordonnée $x$ qui mesure sa distance à ce plan, on peut représenter le potentiel auquel est soumis l'atome d'azote par un puits de potentiel à double minimum, symétrique par rapport à $x=0$. 
\begin{fig}[label]
  \centering
  \begin{subfig}[ammoniac][0.4]
    \includegraphics[width=\columnwidth]{ammoniac.pdf}
  \end{subfig}
  \hspace{1cm}
  \begin{subfig}[ammoniac_model][0.4]
    \includegraphics[width=\columnwidth]{ammoniac_model.pdf}
  \end{subfig}
  \capt{(a) Potentiel de la molécule d'ammoniac en fonction de la position de l'atome d'azote. (b) Modèle de potentiel à double puits rectangulaire.}
\end{fig}
On modélise pour commencer le potentiel de la molécule d'ammoniac par un potentiel à double puits rectangulaire \emph{infinis} de largeur $a$ pour chaque puits, séparés par une barrière de largeur $2b$. On pose donc $V_0=\infty$ sur la \Fig{ammoniac_model}(b). On note $\vphi_n^G(x)$ et $\vphi_n^D(x)$ les modes propres d'ordre $n$ associées aux puits de gauche et de droite, respectivement.
  \begin{quest} 
    \item	Rappeler l'expression de l'énergie $E_n$ associée à ces modes propres, et donner l'expression de la fonction d'onde spatiale $\vphi_n^G(x)$ et $\vphi_n^D(x)$ dans les deux puits.
    \sol{
      L'énergie d'un mode propre d'un puits de potentiel infini de largeur $a$ est donnée par
      \eqbox{
E_n = n^2 \frac{\hbar^2 \pi^2}{2 m a^2}, \quad n=1,2,3,\ldots
      }
      La fonction d'onde spatiale associée à chaque mode propre d'un puits de potentiel infini de bords $x=0$ et $x=a$ s'annule emn dehors de cet intervalle, et est donnée par
      $$
\varphi_n(x) = \sqrt{\frac{2}{a}} \sin\left(\frac{n\pi}{a}x\right), \text{ pour } 0<x<a.
$$
      La fonction d'onde spatiale associée à chaque mode propre s'annule en $x=-b-a$ et $x=-b$ pour le puits de gauche, et en $x=b$ et $x=b+a$ pour le puits de droite. Elles s'obtiennent par translation de la fonction d'onde précédente, soit
      \eqbox{
\vphi_n^G(x) = \sqrt{\frac{2}{a}} \sin\left(\frac{n\pi}{a}(x+b)\right) \quad\text{et}\quad \vphi_n^D(x) = \sqrt{\frac{2}{a}} \sin\left(\frac{n\pi}{a}(x-b)\right).
      }
    }
  \end{quest}
  On considère maintenant que la barrière de potentiel séparant les deux puits est de hauteur finie $V_0$. On cherche les modes propres de plus basse énergie  de ce potentiel à double puits. Pour cela, on construit les fonctions d'ondes symétriques et antisymétriques à partir des modes propres des puits \emph{finis},
  \begin{equation}
\vphi_S(x) = \frac{1}{\sqrt{2}}(\vphi_1^G(x) + \vphi_1^D(x)) \quad\text{et}\quad \vphi_A(x) = \frac{1}{\sqrt{2}}(\vphi_1^G(x) - \vphi_1^D(x)).
  \end{equation}
  Par construction, on a donc $\vphi_S(-x) = \vphi_S(x)$ et $\vphi_A(-x) = -\vphi_A(x)$.
  \begin{quest}
    \item Tracer l'allure des fonctions d'onde $\vphi_S(x)$ et $\vphi_A(x)$ en fonction de $x$ dans l'hypothèse où $V_0$ et $b$ sont \guill{grands}. On reproduira sur le graphe le potentiel à double puits. Que peut-on dire des énergies $E_S$ et $E_A$ associée à ces fonctions d'ondes ?
    \sol{
      Dans l'hypothèse où la barrière est haute ($V_0$ grand) et large ($b$ grand), les fonctions d'ondes $\vphi_S(x)$ et $\vphi_A(x)$ sont très proches de celles des puits infinis, à savoir des fonctions sinusoïdales localisées dans chaque puits.  Dans la barrière, les fonctions d'ondes s'atténuent exponentiellement par effet tunnel.
    \begin{fig}[label]
      \centering
      \includegraphics[width=0.48\columnwidth]{sym_ammoniac.pdf}
      \includegraphics[width=0.48\columnwidth]{antisym_ammoniac.pdf}      
      %\capt{}
    \end{fig}
    Comme les puits sont ici finis, on s'attend à ce que l'énergie $E_S$ du mode symétrique soit inférieure à l'énergie $E_1$ du puits infini simple, elle-même inférieure à l'énergie $E_A$ du mode antisymétrique, soit 
    \eqbox{
E_S < E_1 < E_A.
    }
    }
    \item Résoudre l'équation de Schrödinger indépendante du temps dans les zones I : $-b-a< x < -b$, II : $-b<x<b$ et III : $b<x<b+a$, pour la fonction d'onde  symétrique $\vphi_S(x)$, en tenant compte des conditions de symétrie et de continuité en $x=\pm(b+a)$. On introduira un vecteur d'onde $k_S$ ainsi qu'une distance d'atténuation $\delta$ à exprimer. 
    \sol{
      L'équation de Schrödinger indépendante du temps s'écrit
\begin{equation*}
\begin{cases}
\frac{\d^2\vphi_S}{\d x^2}+\frac{2mE_S}{\hbar^2}\vphi_S=0\quad\text{(zones I et III)},\\[2mm]
\frac{\d^2\vphi_S}{\d x^2}-\frac{2m\left(V_0-E_S\right)}{\hbar^2}\vphi_S=0\quad\text{(zone II)}.
\end{cases}
\end{equation*}
Sachant que la fonction d'onde est symétrique, les solutions s'écrivent
\begin{equation*}
\vphi_S(x)=\begin{cases}
A_S\sin(k_S(x+b)) \text{ si } -b-a< x < -b, \\
B_S\cosh(x/\delta) \text{ si } -b<x<b, \\
A_S\sin(k_S(x-b)) \text{ si } b<x<b+a, \\
\end{cases}
\end{equation*}
où le vecteur d'onde $k_S$ et la distance d'atténuation $\delta$ sont donnés par
\eqbox{
k_S = \frac{\sqrt{2mE_S}}{\hbar} \quad\text{et}\quad \delta = \frac{\hbar}{\sqrt{2m(V_0-E_S)}}.
}
    }
    \item Écrire les conditions de continuité de la fonction d'onde et de sa dérivée en $x=\pm b$. En déduire que le vecteur d'onde $k_S$ vérifie la relation
    \begin{equation}
\tan(ak_S) = -\delta k_S \coth(b / \delta).
    \end{equation}
    Résoudre graphiquement cette équation.

    \sol{
      Comme la fonction d'onde est symétrique, il suffit d'imposer les conditions de continuité en $x=b$. On trouve
\begin{equation*}
\begin{cases}
A_S\sin(k_S a) = B_S\cosh(b/\delta), \\[2mm]\, 
A_S k_S \cos(k_S a) = \frac{B_S}{\delta} \sinh(b/\delta).
\end{cases}
\end{equation*}
En prenant le rapport de ces deux équations, on trouve bien 
% \begin{equation*}
\eqbox{
  \tan(ak_S) = -\delta k_S \coth(b / \delta).
  }
% \end{equation*}
    }
    \item En supposant de plus que la barrière est haute et épaisse, soit $V_0 \gg E_S$, $b \gg \delta$ et $\delta k_S \ll 1$, en déduire une expression approchée de l'énergie $E_S$ du mode symétrique. En reprenant ces calculs pour le mode antisymétrique, montrer que la molécule d'ammoniac oscille entre ces deux positions d'équilibre avec une pulsation d'inversion
    \begin{equation}
\Omega \simeq \frac{\hbar^2 \e^{-\sqrt{8mV_0}\, b /\hbar}}{\sqrt{8m^3V_0} \, a^3}.
    qgs    \end{equation}
    \AN~pour $m=\SI{2.7e-26}{kg}$, $a=\SI{0.1}{nm}$, $b=\SI{0.05}{nm}$ et $V_0=\SI{1}{eV}$.
    \sol{
Dans l'hypothèse où la barrière est haute, la distance d'atténuation est indépendante de l'énergie,
$$
\delta \simeq \frac{\hbar}{\sqrt{2mV_0}},
$$
De plus, on a $\delta k_S \ll 1$.
Dans l'hypothèse où la barrière est épaisse, le rapport $b/\delta$ est grand, et on peut utiliser l'approximation
$$
\coth(b/\delta)  = \frac{1+\e^{-2b/\delta}}{1-\e^{-2b/\delta}} \simeq 1+2\e^{-2b/\delta}.
$$
On trouve alors pour le mode symétrique
$$ 
\tan(ak_S) \simeq -\delta k_S (1+2\e^{-2b/\delta}) \ll 1,
$$
Le mode fondamental vérifie $ak_S = \pi - \varepsilon$ avec $\varepsilon \ll 1$. En développant $\tan(ak_S)$ à l'ordre 1 en $\varepsilon$, on trouve
$$
\varepsilon \simeq \frac{\pi \delta}{a}  (1+2\e^{-2b/\delta}) + o(\delta/a).
$$
soit
$$
E_S  = \frac{\hbar^2 k_S^2}{2m} \simeq \frac{\hbar^2 \pi^2}{2ma^2} \left(1 - \frac{2\varepsilon}{\pi}\right) \simeq \frac{\hbar^2 \pi^2}{2ma^2} - \frac{\hbar^2 \pi^2 \delta}{ma^3}(1+2\e^{-2b/\delta}) \simeq E_{\infty} - E_{\infty} \frac{2\delta}{a} - 2E_{\infty} \frac{\delta}{a} \e^{-2b/\delta},
$$
Le premier terme correspond à l'énergie du mode fondamental d'un puits de potentiel infini de largeur $a$. Le second terme correspond à la correction de l'énergie du mode fondamental d'un puits de potentiel fini de largeur $a$ et de profondeur $V_0$. Le troisième terme correspond à la correction de l'énergie du mode fondamental en raison du couplage entre les deux puits. Ce terme de couplage change de signe pour le mode antisymétrique, de sorte que l'écart d'énergie entre les modes symétrique et antisymétrique est donné par
$$
E_A - E_S \simeq 4E_{\infty} \frac{\delta}{a} \e^{-2b/\delta} = \frac{2\hbar^2 \pi^2 \delta}{ma^3} \e^{-2b/\delta} = \frac{2\hbar^2 \pi^2}{ma^3} \frac{\hbar}{\sqrt{2mV_0}} \e^{-2b\sqrt{2mV_0}/\hbar}.
$$
Avec $E_A - E_S = \hbar \Omega$, on trouve
\eqbox{
\Omega \simeq \frac{\hbar^2 \e^{-\sqrt{8mV_0}\, b /\hbar}}{\sqrt{8m^3V_0} \, a^3}.
}

    }
  \end{quest}
\end{exo}

%%%%%%%%%%%%%%%%%%%%%%%%%%%%%%%%%%%
\begin{exo}[PC][2][colle]{Marche de potentiel}
  Une particule de masse $m$ et d'énergie $E$ arrive de $x=-\infty$ sur une marche de potentiel de hauteur $V_0>E$ située en $x=0$.
  \begin{quest}
  \item Que se passe-t-il en mécanique classique?
\sol{
La particule rebondit sur la barrière.
}
  \item On s'intéresse à des états stationnaires, de fonction d'onde $\psi(x, t)=\varphi(x) \exp \left(\frac{-i E}{\hbar} t\right)$. Déterminer $\varphi(x)$ dans les différents domaines. On posera $k=\sqrt{\frac{2 m E}{\hbar^2}}$ et $\alpha=\sqrt{\frac{2 m\left(V_0-E\right)}{\hbar^2}}$.
  \sol{

  \udl{Équation de Schrödinger} indépendante du temps s’écrit ici
  \begin{equation*}
  \begin{cases}
  \frac{\d^2\varphi}{\d x^2}+\frac{2mE}{\hbar^2}\varphi=0\quad\text{(région I, x<0)},\\[2mm]
  \frac{\d^2\varphi}{\d x^2}-\frac{2m\left(V_0-E\right)}{\hbar^2}\varphi=0\quad\text{(région II, x>0)}.
  \end{cases}
  \end{equation*}
  En posant  $k=\sqrt{\frac{2mE}{\hbar^2}}$ et $\alpha=\sqrt{\frac{2m\left(V_0-E\right)}{\hbar^2}}$, la partie spatiale de la fonction d'onde s'écrit :
  \begin{equation*}
  \varphi(x)=\begin{cases}
  A_\mathrm{I}\,\e^{ikx}+B_\mathrm{I}\,\e^{-ikx} \text{ si } x<0, \\
  A_\mathrm{II}\e^{\alpha x}+B_\mathrm{II}\e^{-\alpha x} \text{ si } 0<x, \\
  \end{cases}
  \end{equation*}

  \udl{Conditions aux limites} : Pour éviter une divergence en $x\rightarrow +\infty$, on doit imposer $A_\mathrm{II}=0$.
  On impose ensuite la continuité de $\vphi$ et $\dd{\vphi}{x}$ en $x=0$.

  La fonction d'onde a donc la forme
  \begin{equation*}
  \varphi(x)=\begin{cases}
  A_\mathrm{I}\,\e^{ikx}+B_\mathrm{I}\,\e^{-ikx} \text{ si } x<0, \\
  B_\mathrm{II}\e^{-\alpha x} \text{ si } 0<x, \\
  \end{cases}
  \end{equation*}
  où $A_\mathrm{I}$ est l'amplitude de l'onde incidente et $B_\mathrm{I}$ l'amplitude de l'onde réfléchie. L'onde est évanescente pour $x>0$. Les conditions aux limites donnent
  \begin{equation*}
  \begin{cases}
  A_\mathrm{I}+B_\mathrm{I} = B_\mathrm{II} \\
  ik (A_\mathrm{I}-B_\mathrm{I}) =  \alpha B_\mathrm{II} =  \alpha(A_\mathrm{I}+B_\mathrm{I}) \\
  \end{cases} \Rar
  \begin{cases}
  B_\mathrm{II} = A_\mathrm{I}+B_\mathrm{I} =  \frac{2ik }{\alpha+ik}  A_\mathrm{I}  \\
  B_\mathrm{I} =  -\frac{\alpha-ik }{\alpha+ik}  A_\mathrm{I} \\
  \end{cases}
  \end{equation*}
  }
  \item Définir et calculer le coefficient de réflexion en amplitude $\underline{r}$ puis son module.
  \sol{
On définit
$$
\udl{r} = \frac{B_\mathrm{I}}{A_\mathrm{I}}  =  -\frac{\alpha-ik }{\alpha+ik},
$$
de module $|\udl{r}| = 1$.
  }
  \item On donne la densité de courant de probabilité $j(x, t)=\frac{i \hbar}{2 m}\left(\psi \frac{\partial \psi^*}{\partial x}-\psi^* \frac{\partial \psi}{\partial x}\right)$. Donner les vecteurs densité de courant de probabilité incident et réfléchi $j_i$ et $j_r$. Que peut-on dire du coefficient de réflexion en probabilité défini comme $R=\left|\frac{j_r}{j_i}\right|$ ?
  \sol{

  La densité de courant de probabilité incidente vaut
  $$
j_i = -\frac{\hbar}{m}\im\left(\psi_i \frac{\partial \psi_i^*}{\partial x}\right),
  $$
  avec $\psi_i = A_\mathrm{I}\exp \left(\frac{-i E}{\hbar} t\right) \e^{ikx}$, soit
  $$
j_i = \frac{\hbar k}{m}.
  $$
  De même pour la densité de courant de probabilité réfléchie,
  $$
j_r = -\frac{\hbar k}{m} |\udl{r}|^2.
  $$
  On en déduit le coefficient de réflexion
  $$
  R = |\udl{r}|^2 = 1.
  $$
  Il y a réflexion totale : l'onde est entièrement réfléchie à la barrière.

  \udl{Remarque} : Comme $-\frac{\alpha-ik }{\alpha+ik}$ est de module $1$, on peut écrire
  $$
  -\frac{\alpha-ik }{\alpha+ik} = e^{i(\pi-2\theta)}, \quad \text{avec}\quad \tan \theta = \frac{k}{\alpha}.
  $$
  L'onde réfléchie accumule un certain retard $\tau = \frac{\pi-2\theta}{\omega}$ avec $E=\hbar \omega$ parce qu'elle explore une partie de la région II interdite.
  }
  \item Déterminer le facteur de transmission de la barrière $T$.
  \sol{
  $T = 1-R=0$.
  }
\end{quest}
\end{exo}
%%%%%%%%%%%%%%%%%%%%%%%%%%%%%%%%%%%


%%%%%%%%%%%%%%%%%%%%%%%%%%%%%%%%%%%
\begin{exo}[PC][1][colle]{Résonance}
  On étudie une particule libre venant de $x\to-\infty$ qui rencontre en $x=0$ une barrière de potentiel d'épaisseur $a$ et de hauteur $V_0>0$. On cherche des états stationnaires d’énergie $E>V_0$.
  % est supposée supérieure à la a
  % Autrement dit, la particule évolue dans le potentiel schématisé Fig.~\ref{fig:barriere} tel que
  % \begin{equation}
  % V(x)=\begin{cases}
  % 0 \text { si } x<-\frac{a}{2},\\
  % V_0 \text{ si } -\frac{a}{2}<x<\frac{a}{2},\\
  % 0 \text { si } x>\frac{a}{2}.
  % \end{cases}
  % \end{equation}
  %
  % \begin{fig}[]
  %   \centering
  %   \includegraphics[width=0.4\textwidth]{barriere.pdf}
  %   \capt{Représentation du potentiel $V(x)$ dans le cas de la barrière de potentiel de hauteur $V_0$.}
  % \end{fig}

\begin{quest}
  \item  Écrire l’équation de Schrödinger dans chaque région, et en déduire la forme de la fonction d’onde spatiale $\vphi(x)$ dans chaque région.
  \item Déterminer les constantes à partir des conditions aux limites.
  \item En déduire l’expression du coefficient de  transmission $T$. Tracer $T$ en fonction de l’épaisseur $a$ de la barrière et Interpréter.
\end{quest}
\end{exo}
%%%%%%%%%%%%%%%%%%%%%%%%%%%%%%%%%%%

%%%%%%%%%%%%%%%%%%%%%%%%%%%%%%%%%%%%%%%%%%%%%%%%%%%%%%%%%%%%%%%%%%
\begin{exo}[PC][1][TD]{La marche de potentiel}
  \Notions{\item Équation de Schrödinger \task Onde stationnaire \task Conditions aux limites}
On étudie le cas d'une particule qui rencontre une marche de potentiel représentée par
\begin{equation*}
V(x) =
\begin{cases}
0 \mbox{ si } x<0, \\
V_0 \mbox{ si } x> 0,
\end{cases}
\end{equation*}
avec $V_0>0$. On souhaite étudier l'existence d'états stationnaires d'énergie $E$ pour lesquels $0<E<V_0$.
\begin{quest}
\item Écrire l'équation de Schrödinger indépendante du temps sur la fonction d'onde spatiale $\vphi(x)$ dans chacune des régions $x<0$ et $x>0$, puis donner la forme générale de la solution dans chaque région.
\sol{
Cette situation correspond à un puits de potentiel de hauteur finie, mais qui occuperait tout le demi-espace $x>0$. L'\udl{équation de Schrödinger} indépendante du temps s’écrit ici
% \udl{Équation de Schrödinger} indépendante du temps s’écrit ici
\begin{equation*}
\begin{cases}
\frac{\d^2\varphi_I}{\d x^2}+\frac{2mE}{\hbar^2}\varphi_I=0\quad\text{(région I, x<0)},\\[2mm]
\frac{\d^2\varphi_{II}}{\d x^2}-\frac{2m\left(V_0-E\right)}{\hbar^2}\varphi_{II}=0\quad\text{(région II, x>0)},
\end{cases}
\end{equation*}
en notant que $V_0-E>0$.
On peut poser $k=\sqrt{\frac{2mE}{\hbar^2}}$ le vecteur d'onde dans la région I et $q=\sqrt{\frac{2m\left(V_0-E\right)}{\hbar^2}}$ la partie imaginaire du vecteur d'onde dans la région II. La partie spatiale de la fonction d'onde s'écrit
% \fbox{
  \begin{center}
    \boxed{
      $$
      \varphi(x)=\begin{cases}
        A_\mathrm{I}\,\e^{ik_\mathrm{I}x}+B_\mathrm{I}\,\e^{-ik_\mathrm{I}x} \text{ si } x<0, \\
        A_\mathrm{II}\e^{qx}+B_\mathrm{II}\e^{-qx} \text{ si } x>0. \\
      \end{cases}
      $$
      }
  \end{center}
% }



% En particulier, on note que la relation $R+T=1$ qui traduit la (associé à la conservation de la norme de la fonction d'onde) est bien vérifiée.

}
\item Simplifier la forme de la solution en imposant que $\vphi(x)$ ne diverge pas. Définir dans la région $x<0$ une onde incidente et une onde réfléchie, et discuter de la nature de l'onde transmise dans la région $x>0$. Écrire les conditions aux limites en $x=0$ et en déduire le coefficient de réflexion $R$.
\sol{
  \udl{Conditions aux limites} : Pour éviter une divergence en $x\rightarrow +\infty$, on doit imposer $A_\mathrm{II}=0$.
La fonction d'onde a donc la forme
\begin{equation*}
\varphi(x)=\begin{cases}
A_\mathrm{I}\,\e^{ikx}+B_\mathrm{I}\,\e^{-ikx} \text{ si } x<0, \\
B_\mathrm{II}\e^{- qx} \text{ si } 0<x, \\
\end{cases}
\end{equation*}
où $A_\mathrm{I}$ est l'amplitude de l'onde incidente et $B_\mathrm{I}$ l'amplitude de l'onde réfléchie. L'onde est évanescente pour $x>0$. Les conditions aux limites en $x=0$ donnent ensuite
\begin{equation*}
\begin{cases}
A_\mathrm{I}+B_\mathrm{I} = B_\mathrm{II} \\
ik (A_\mathrm{I}-B_\mathrm{I}) =  q B_\mathrm{II} =  q(A_\mathrm{I}+B_\mathrm{I}) \\
\end{cases} \Rar
\begin{cases}
B_\mathrm{II} = A_\mathrm{I}+B_\mathrm{I} =  \frac{2ik }{q+ik}  A_\mathrm{I}  \\
B_\mathrm{I} =  -\frac{q-ik }{q+ik}  A_\mathrm{I} \\
\end{cases}
\end{equation*}
On en déduit le coefficient de réflexion en amplitude $r= \frac{B_\mathrm{I}}{A_\mathrm{I}} = -\frac{q-ik }{q+ik}$ puis en probabilité
$$
R =|r|^2 = \left| \frac{q-ik }{q+ik} \right|^2 = 1.
$$
Il y a \fbox{réflexion totale} : l'onde est entièrement réfléchie par la barrière.
}
\item Déterminer la longueur de pénétration $\delta$ de l'onde dans la marche de potentiel ainsi que le retard $\tau$ que l'onde accumule lors de la réflexion. \AN~pour un électron de dans un métal d'énergie $E=\SI{1}{eV}$ qui rencontre le vide, de potentiel $V_0=\SI{4}{eV}$.
\sol{
  Comme $r$ est de module 1 , on peut écrire
  $$
  r=e^{i(\pi-2 \theta)}, \quad \text { avec } \quad \tan \theta=k \delta,
  $$
  où $\delta=1 / q$ est la longueur de pénétration de l'onde dans la barrière. L'onde réfléchie accumule un certain retard $\tau$ tel que $\omega \tau=\pi-2 \theta$ et $E=\hbar \omega$, parce qu'elle explore une partie de la région II interdite. \AN on trouve \fbox{$\delta = \SI{}{}$} et \fbox{$\tau = \SI{}{}$}
}
\end{quest}
\end{exo}
%%%%%%%%%%%%%%%%%%%%%%%%%%%%%%%%%%%%%%%%%%%%%%%%%%%%%%%%%%%%

%%%%%%%%%%%%%%%%%%%%%%%%%%%%%%%%%%%%%%%%%%%%%%%%%%%%%%%%%%%%
\begin{exo}[PC][3][oral]{Modèle de Gamow de la radioactivité $\alpha$}
  \Notions{\item Équation de Schrödinger \task Onde stationnaire \task Effet tunnel}
  La radioactivité $\alpha$ est un type de radioactivité où un noyau instable de numéro atomique $Z$ se désintègre en un noyau plus petit de numéro atomique $Z-2$ en émettant un atome d'hélium, aussi appelé particule~$\alpha$. 
  Les données expérimentales représentées ci-dessous, vérifient la loi de Geiger-Nutal, qui donne la lien entre la durée de demie-vie $t_{1/2}$ de l'atome et l'énergie $Q_\alpha$ de la particule $\alpha$ émise, 
  \begin{equation}
  \ln t_{1/2} = \frac{A}{\sqrt{Q_\alpha}} + B,
  \end{equation}
  où $A$ et $B$ sont deux constantes qui dépendent de la nature du noyau émetteur.

  \begin{fig}[geiger_nutal]
    \centering
    \includegraphics[width=0.7\columnwidth]{geiger_nutal.png}
    %\capt{}
  \end{fig}

  G. \textsc{Gamow} a proposé un modèle simplifié qui repose sur l'effet tunnel pour décrire la radioactivité $\alpha$. Lorsque la particule $\alpha$ est à l'intérieur du noyau de rayon $R$, elle ressent le potentiel d'interaction forte qui assure la cohésion du noyau, modélisé par un puits de profondeur $V_0$ et de largeur $2R$. Lorsque la particule $\alpha$ sort du noyau, elle est soumise au potentiel de répulsion électrostatique $V(x)$ généré par les $Z-2$ protons restants dans le noyau. La particule $\alpha$ rebondit périodiquement dans le puits de potentiel. \\ 
  
  Retrouver la loi de Geiger-Nutal à partir du modèle de Gamow.

% En précisant les hypothèses et en expliquant la démarche, démontrer la loi de Geiger et Nuttal,
%   \begin{equation}
%   \ln\tau_{1/2}=\frac{A}{\sqrt{E}}+B,
%   \label{eqn:geiger_nuttal}
%   \end{equation}
%   où $\tau_{1/2}$ représente le temps de demi-vie du noyau radioactif, $E$ est l'énergie de la particule $\alpha$, et $A$ et $B$ sont deux constantes.
  %
  \sol{

  \begin{fig}[V_alpha]
    \centering
    \includegraphics[width=0.5\textwidth]{alpha.pdf}
    \capt{Illustration du potentiel $V(x)$ auquel est soumis une particule $\alpha$ dans le modèle de Gamow. Les rectangles bleus représentent une approximation de la barrière par une succession de barrières rectangulaires.} %On peut donner quelques ordres de grandeur : $V_0=(25-40)\SI{}{MeV}$, $R=\SI{5}{fm}$, $E=(4-9)\SI{}{MeV}$.Modèle de Gamow de la radioactivité $\alpha$}
  \end{fig}
  }

\end{exo}
%%%%%%%%%%%%%%%%%%%%%%%%%%%%%%%%%%%

%%%%%%%%%%%%%%%%%%%%%%%%%%%%%%%%%%%%%%%%%%%%%%%%%%%%%%%%%%%%%%%%%%%%%%
\begin{exo}[PC][2][TD]{Effet Ramsauer}
  \Notions{\item Équation de Schrödinger \task Onde stationnaire \task État libre \task Transmission}
En 1921, Ramsauer avait constaté que pour certaines valeurs particulières de l'énergie incidente, des gaz rares (hélium, argon, néon ...) étaient parfaitement transparents à des faisceaux d'électrons de basse énergie. Ce phénomène peut se comprendre dans le modèle unidimensionnel suivant. On considère une solution stationnaire de l'équation de Schrödinger d'énergie positive $E>0$ pour une particule de masse $m$ dans un puits de potentiel fini de profondeur $V_0>0$, soit
\begin{equation*}
V(x) =
\begin{cases}
0 \mbox{ si } x<-a, \\
-V_0 \mbox{ si } -a\leq x\leq a,\\
0 \mbox{ si } x>a. \\
\end{cases}
\end{equation*}

\begin{quest}
  \item On pose $q=\sqrt{\frac{2m(V_0+E)}{\hbar^2}}$ et $k=\sqrt{\frac{2mE}{\hbar^2}}$. Justifier qu'on cherche une solution de la forme
  \begin{equation*}
\varphi(x)=
\begin{cases}
\e^{ikx}+A\,\e^{-ikx} \mbox{ si } x\leq -a, \\
B\,\e^{iqx}+C\,\e^{-iqx}  \mbox{ si } -a\leq x\leq a,\\
D\,\e^{ikx}  \mbox{ si } x\geq a. \\
\end{cases}
\end{equation*}
\sol{
Comme $E>V$ dans les trois régions, la fonction d'onde s'écrit comme une OPPH complexe. Les vecteurs d'onde associés sont $k$ dans les régions I et III, de potentiel nul, et $q$ dans la région II, de potentiel $-V_0$. Il n'y a pas d'onde provenant de $+\infty$, donc pas de terme $e^{-ikx}$ pour $x\geq a$. Enfin, on a choisi d'attribuer le coefficient $1$ à l'onde incidente plutôt que de normaliser la fonction d'onde ; cela simplifie les calculs et ne pose pas de problème si l'on se contente de calculer des coefficients de réflexion ou transmission.
}
\item Écrire les relations de continuité vérifiées en $x=\pm a$.
% \item En posant $\Delta=(q+k)^2-4\,\e^{4iqa}(q-k)^2$, on peut montrer que le système précédent admet pour solution $A=\frac{k^2-q^2}{\Delta}\,\e^{-2ika}(1-\e^{4iqa})$ et $D=\frac{4kq}{\Delta}\e^{-2i(k-q)a}$. Calculer la probabilité de transmission $T=\abs{D}^2$ puis de réflexion $R=\abs{A}^2$, et vérifier explicitement que $R+T=1$.
Définir le coefficient de transmission de probabilité $T$ puis l'exprimer en fonction des coefficients $A,B,C,D$.
\end{quest}
On définit $\mu = \frac{V_0^2}{4E(V_0+E)}$. 
On peut montrer que le coefficient de transmission a pour expression
$$
T = \frac{1}{1+\mu \sin^2(2qa)},
$$

\begin{quest}
  \item 
  Que vaut le coefficient de réflexion $R$ ?
\sol{
Continuité de $\vphi$ et $\vphi'$ en $x=\pm a$. Donc
$$
\begin{cases}
1+A = B + C \\
ik(1-A) = iq(B-C)\\
\end{cases}
\qquad
\begin{cases}
Be^{iqa} + Ce^{-iqa} = De^{ika} \\
iq(Be^{iqa} - Ce^{-iqa}) = ikDe^{ika}\\
\end{cases}
$$
Par définition, $T = \frac{J_t}{J_i}$ où $J =\frac{\hbar K}{m} \vphi \vphi^\star$ est la densité de courant de probabilité. Or $\vphi_i(x) = e^{ikx}$ et $\vphi_t(x) = D e^{ikx}$ donc ici $T = |D|^2$.
}
\item Montrer qu'il existe certaines valeurs de l'énergie pour lesquelles $T=1$. Interpréter ce résultat.
\sol{
On simplifie $\mu = \frac{4m^2 V_0^2}{4\hbar^4} \frac{\hbar^2}{2mE} \frac{\hbar^2}{2m(V_0+E)} = \frac{V_0^2}{4E(V_0+E)}$. Ce paramètre dépend de $E$ mais c'est surtout la dépendance de $\sin^2(2qa)$ qui est importante. En effet, ce terme s'annule si $2qa = n \pi$ avec $n\geq0$ entier, soit $E = n^2 \frac{\hbar^2\pi^2}{8ma^2}-V_0 = n^2 E_1-V_0$
avec $E_1=\frac{\hbar^2\pi^2}{8ma^2}$ l'énergie du mode fondamental du puits infini de largeur $2a$. La condition de quantification du vecteur d'onde $q$ et de l'énergie $E$ est la même que pour le puits infini !
}
\item Pour l'hélium, l'énergie incidente à laquelle apparaît le phénomène est $E=\SI{0,7}{eV}$. En supposant l'atome d'hélium de rayon $a=\SI{0,1}{nm}$, calculer la profondeur  du puits $V_0$ à l'intérieur de l'atome.
\sol{
$E=\frac{\hbar^2\pi^2}{2ma^2}-V_0 = \SI{0,7}{eV}$, d'où $V_0 = \frac{\hbar^2\pi^2}{8ma^2}-E = \SI{8.7}{eV}$ avec $m=\SI{9.1e-31}{kg}$ la masse de l'électron.
}
\item Comment se comporte le coefficient de transmission $T$ lorsque $E\to 0$ ? Lorsqu'on envoie des atomes d'hydrogène très lents sur une surface d'hélium liquide, on constate que ces atomes rebondissent élastiquement au lieu d'être adsorbés. Interpréter qualitativement ce phénomène.
\sol{
On a alors $\sin^2(2qa) \rightarrow \sin^2(2q_0a) \neq 0$ a priori, avec $q_0 = \sqrt{\frac{2mV_0}{\hbar^2}}$ et $\mu \sim \frac{V_0}{4E}$, soit
$$
T \sim \frac{4E}{V_0 \sin^2(2q_0a)} \rightarrow 0.
$$
Les particules sont toutes réfléchies.
}
\end{quest}
\end{exo}
%%%%%%%%%%%%%%%%%%%%%%%%%%%%%%%%%%%%%%%%%%%%%%%%%%%%%%%%%%%%%%%%%%

%%%%%%%%%%%%%%%%%%%%%%%%%%%%%%%%%%%
\begin{exo}[PC][2][TD]{Particule piégée par un pic de potentiel}
  \Notions{\item Équation de Schrödinger \task Onde stationnaire \task Confinement}
  On considère une particule de masse $m$ dans un potentiel nul partout, sauf en $x=0$, qui vérifie
  \begin{equation*}
    V(x) =
    \begin{cases}
    0 \mbox{ si } x<0, \\
    0 \mbox{ si } x> 0,
    \end{cases} \quad \text{ et } \int_{-\veps}^{+\veps} V(x) \d x = -V_0 b,~\forall \veps > 0.
    \end{equation*}
  Ce modèle peut représenter un électron piégé par un atome ou un défaut ponctuel dans un cristal, dont le puits de potentiel est de profondeur $V_0>0$ infiniment grande et de largeur $b>0$ infiniment petite. On suppose que la particule est dans un état lié d'énergie $E<0$.
  \begin{quest}
    \item	On cherche de solutions stationnaires de l'équation de Schrödinger de la forme $\psi(x,t) = \varphi(x) \e^{-iEt/\hbar}$. Écrire l'équation de Schrödinger indépendante du temps qui porte sur $\vphi(x)$ dans les régions I ($x<0$) et II ($x>0$) puis trouver la forme des solutions.
    \sol{
    L'équation de Schrödinger indépendante du temps est la même dans les deux régions, soit 
    \begin{equation*}
    % \begin{cases}
    \frac{\d^2\varphi}{\d x^2}+\frac{2mE}{\hbar^2}\varphi= 0 \Rar \vphi''(x) - q^2\vphi(x) = 0,
    % \quad\text{(région I, 0<x<0)},\\[2mm]
    % \frac{\d^2\varphi}{\d x^2}=0\quad\text{(région II, x>0)}.
    % \end{cases}
    \end{equation*}
    où l'on a posé $q=\sqrt{\frac{2m|E|}{\hbar^2}}$. La solution a la même forme dans chaque région mais avec des coefficients différents a priori,
    \begin{equation*}
    \begin{cases}
      \varphi_{I}(x)=A = \,\e^{qx}+B\,\e^{ -qx} \text{ si } x<0, \\
      \varphi_{II}(x)=C\,\e^{qx}+D\,\e^{-qx} \text{ si } x>0. \\
    \end{cases}
    \end{equation*}
    }
    
    \item Quelles conditions aux limites portent sur la fonction d'onde spatiale $\vphi(x)$ ? Sur sa dérivée $\vphi'(x)$ ? En déduire trois des quatre coefficients d'intégration $A,B,C,D$.
    \sol{
      La fonction d'onde ne diverge pas en $x\rightarrow \pm\infty$ donc $B=0$ et $C=0$. 
    La fonction d'onde $\vphi(x)$ est continue en $x$ mais pas sa dérivée $\vphi'(x)$ car le potentiel subit une discontinuité infini en $x=0$. On a alors 
    $$
    \vphi_{I}(x=0) = \vphi_{II}(x=0) \Rar A=D.
    $$
    On en déduit que la fonction d'onde est de la forme
    % \begin{equation*}
      \eqbox{
        \varphi(x)=A\,\e^{ -q|x|}.
      }
    % \end{equation*}
    }
    \item Déterminer et représenter graphiquement la densité de probabilité de présence $\rho(x)$. Déterminer le dernier coefficient à partir de la condition de normalisation. 
    \sol{
      La densité de probabilité de présence est donnée par $\rho(x) = |\vphi(x)|^2 = A^2\e^{-2q|x|}$. La condition de normalisation impose
      La normalisation impose 
      $$
      \int_{-\infty}^{+\infty} \rho(x) \d x = 1 \Rar 2A^2 \int_{0}^{+\infty} \e^{-2qx} \d x = 1 \Rar A = \sqrt{\frac{q}{2}}.
      $$
      On en déduit 
      \eqbox{
        \varphi(x)=\sqrt{\frac{q}{2}}\,\e^{ -q|x|}.
      }
      }
      \item La valeur de l'énergie $E$ est contrainte par le potentiel. Pour le montrer, écrire l'équation de Schrödinger indépendante du temps pour une énergie $E$ pour $x$ quelconque. Intégrer cette équation sur un petit intervalle $[-\veps,+\veps]$ et étudier la limite $\veps \xrightarrow{} 0$. En déduire que l'énergie vaut 
        $$E = -\frac{2mV_0^2b^2}{\hbar^2}.$$
      %  et montrer que l'on retrouve la condition de normalisation.
      \sol{
      L'équation de Schrödinger indépendante du temps est
      $$
      \frac{\d^2\varphi}{\d x^2} - \left(q^2+\frac{2m}{\hbar^2}V(x)\right)\varphi = 0.
      $$
      En intégrant cette équation sur un intervalle $[-\veps,+\veps]$, on obtient
      $$
      \int_{-\veps}^{+\veps} \left[\vphi''(x) - \left(q^2+\frac{2m}{\hbar^2}V(x)\right)\varphi(x)\right] \d x = 0 \Rar \int_{-\veps}^{+\veps} \vphi''(x) \d x - q^2\int_{-\veps}^{+\veps} \varphi(x) \d x + \frac{2m}{\hbar^2} \int_{-\veps}^{+\veps} V(x) \vphi(x) \d x = 0.
      $$ 
      Dans la limite $\veps \xrightarrow 0$, on trouve 
      $$
      \vphi_{II}'(0)-\vphi_{I}'(0) - \vphi(0) \frac{2mV_0b}{\hbar^2} = 0 \Rar -qA - \frac{2mV_0b}{\hbar^2}A = 0 \Rar q = \frac{2mV_0b}{\hbar^2},
      $$
      et donc 
      $$
      E = -\frac{\hbar^2q^2}{2m} = -\frac{\hbar^2}{2m}\left(\frac{2mV_0b}{\hbar^2}\right)^2,
      $$
      soit 
      \eqbox{
        E = -\frac{2mV_0^2b^2}{\hbar^2}.
      }
      }
  \end{quest} 
\end{exo}
%%%%%%%%%%%%%%%%%%%%%%%%%%%%%%%%%%%

%%%%%%%%%%%%%%%%%%%%%%%%%%%%%%%%%%%
\begin{exo}[PC][3][TD]{Modèle de Kronig-Penney}
  \Notions{\item Équation de Schrödinger \task Onde stationnaire \task Puits de potentiel périodique}
  On étudie le comportement d'un électron de conduction dans un matériau cristallin. On modélise le cristal par une chaîne de cations espacés régulièrement d'une distance $a$ suivant l'axe $x$. L'électron subit le potentiel $V(x)$ produit par la chaîne de cations, qui est donc périodique de période $a$, soit 
  $$
  V(x+a) = x,~\forall x \in \R.
  $$
  L'allure réelle du potentiel est représentée à gauche ci-dessous. Nous modélisons ce potentiel par une suite de pics de potentiel de profondeur $V_0>0$ infiniment grande et de largeur $b>0$ infiniment petite, répartis tous les $x_n = n a$ avec $n\in \Z$, comme illustré sur la figure de droite. On a donc 
  \begin{equation*}
    V(0 < x < a) = 0 \quad \text{ et } \int_{-\veps}^{+\veps} V(x) \d x = -V_0 b,~\forall \veps > 0.
    \end{equation*}
    \begin{fig}[label]
      \centering
      \begin{subfig}[label][0.47]
        \includegraphics[width=\columnwidth]{potentiel_reseau.pdf}
      \end{subfig}
      \begin{subfig}[label][0.47]
        \includegraphics[width=\columnwidth]{potentiel_reseau_pics.pdf}
      \end{subfig}
      \capt{
        (a) Potentiel réeliste produit par une chaîne de cations d'un cristal séparés par la distance inter-atomique $a$. Le puits de potentiel produit par un seul cation est de profondeur $V_0$ et de largeur typique $b$.
        (b) Modélisation du potentiel par un peigne de pics de potentiel.
      }
    \end{fig}
    On cherche un état libre pour l'électron d'énergie $E>0$ et de vecteur d'onde effectif $\alpha=\sqrt{\frac{2mE}{\hbar^2}}$.
    \begin{quest}
      \item	On cherche une solution stationnaire de l'équation de Schrödinger de la forme $\psi(x,t) = \varphi(x) \e^{-iEt/\hbar}$. Écrire l'équation de Schrödinger indépendante du temps qui porte sur $\vphi(x)$ dans la région $0<x<a$, puis trouver la forme de la solution en fonction de $\alpha$ et de  deux coefficients d'intégration $A$ et $B$. 
      \sol{
      L'équation de Schrödinger indépendante du temps est la même dans les deux régions, soit 
      \begin{equation*}
      % \begin{cases}
      \frac{\d^2\varphi}{\d x^2}+\frac{2mE}{\hbar^2}\varphi= 0 \Rar \vphi''(x) + \alpha^2\vphi(x) = 0,
      % \quad\text{(région I, 0<x<0)},\\[2mm]
      % \frac{\d^2\varphi}{\d x^2}=0\quad\text{(région II, x>0)}.
      % \end{cases}
      \end{equation*}
      où l'on a posé $\alpha=\sqrt{\frac{2mE}{\hbar^2}}$. La solution a la même forme dans chaque région mais avec des coefficients différents a priori,
      \begin{equation*}
      % \begin{cases}
        \varphi(0 <x < a) = A \,\e^{qx}+B\,\e^{ -qx} \text{ si } x<0, 
        % \\
        % \varphi_{II}(x)=C\,\e^{qx}+D\,\e^{-qx} \text{ si } x>0. \\
      % \end{cases}
      \end{equation*}
      }
      \item Quelle condition aux limites porte sur la fonction d'onde $\vphi(x)$ en $x=0$ ? Le théorème de Bloch affirme que dans un potentiel de période $a$ la fonction d'onde spatiale peut s'écrire sous la forme $$\varphi(x) = u(x)\, \e^{ikx},$$
       où $u(x)$ est une fonction périodique de période $a$. La grandeur $k$ est appelé moment cristallin. Justifier cette forme de la fonction d'onde en invoquant la densité de probabilité. Quelle autre condition peut-on alors écrire sur $\vphi(x)$ ?
       \item À partir des conditions précédentes, et sachant que les coefficients $A$ et $B$ ne peuvent pas être nuls, déterminer une équation implicite qui relie $k$ et $\alpha$. La résoudre graphiquement et tracer le graphe de la fonction $E(k)$, puis interpréter comme l'existence de bandes d'énergie permises et interdites. Comme se simplifie la fonction $E(k)$ aux faibles moments cristallins $ka \ll 1$ ? Commenter.
    \end{quest}
\end{exo}
%%%%%%%%%%%%%%%%%%%%%%%%%%%%%%%%%%%


%%%%%%%%%%%%%%%%%%%%%%%%%%%%%%%%%%%%%%%%%%%%%%%%%%%%%%%%%%%%%%%%%%%%%%%%%%%%%
\begin{exo}[PC][2][TD]{Puits dissymétrique}
  \Notions{\item Équation de Schrödinger \task Onde stationnaire \task Confinement}
On étudie les états stationnaires d'une particule liée d'énergie $E$ telle que $-V_0<E<0$ dans un puits de potentiel dissymétrique de la forme
\begin{equation*}
V(x)=\begin{cases}
\infty \mathrm{~~si~~} x<0,\\
-V_0 \mathrm{~~si~~} 0<x<L,\\
0 \mathrm{~~si~~} x>L.
\end{cases}
\end{equation*}
\begin{quest}
\item Déterminer la forme de la fonction d'onde spatiale dans chaque région, sans déterminer pour l'instant la valeur des 3 constantes d'intégration non nulles.
\sol{

\udl{Équation de Schrödinger} indépendante du temps s’écrit ici
\begin{equation*}
\begin{cases}
\frac{\d^2\varphi}{\d x^2}+\frac{2m(V_0+E)}{\hbar^2}\varphi=0\quad\text{(région I, 0<x<L)},\\[2mm]
\frac{\d^2\varphi}{\d x^2}-\frac{2m |E|}{\hbar^2}\varphi=0\quad\text{(région II, x>L)}.
\end{cases}
\end{equation*}
car $E<0$. En posant $k_\mathrm{II}=\sqrt{\frac{2m|E|}{\hbar^2}}$ et $k_\mathrm{I}=\sqrt{\frac{2m\left(V_0+E\right)}{\hbar^2}}$, la partie spatiale de la fonction d'onde s'écrit :
\begin{equation*}
\varphi(x)=\begin{cases}
A_\mathrm{I}\,\e^{ik_\mathrm{I}x}+B_\mathrm{I}\,\e^{-ik_\mathrm{I}x} \text{ si } L>x>0, \\
A_\mathrm{II}\e^{k_\mathrm{II}(x-L)}+B_\mathrm{II}\e^{-k_\mathrm{II}(x-L)} \text{ si } x>L, \\
\end{cases}
\end{equation*}

\udl{Conditions aux limites} : Pour éviter une divergence en $x\rightarrow +\infty$, on doit imposer $A_\mathrm{II}=0$.
On impose ensuite la continuité de $\vphi$ en $x=0$ et $x=L$ et de $\dd{\vphi}{x}$ en $x=L$, avec $\vphi(x)=0$ pour $x<0$. Attention : le potentiel est infini pour $x<0$, et donc $\dd{\vphi}{x}$ n'est pas continu en ce point.

}
\item À partir des conditions aux limites, déterminer les relations satisfaites par les constantes d'intégration.
\sol{
La fonction d'onde a donc la forme
\begin{equation*}
\varphi(x)=\begin{cases}
A_\mathrm{I}\,\cos(k_\mathrm{I}x)+B_\mathrm{I}\,\sin(k_\mathrm{I}x) \text{ si } L>x>0, \\
B_\mathrm{II}\e^{-k_\mathrm{II}(x-L)} \text{ si } x>L. \\
\end{cases}
\end{equation*}
 L'onde est évanescente pour $x>0$. Les conditions aux limites  en $x=0$ donnent $A_\mathrm{I}=0$. Puis en $x=L$, on a
 \begin{equation*}
 \begin{cases}
 B_\mathrm{I} \sin(k_\mathrm{I}L) = B_\mathrm{II} \\
 k_\mathrm{I} B_\mathrm{I} \cos(k_\mathrm{I}L) = - k_\mathrm{II} B_\mathrm{II}.
 \end{cases}
 \end{equation*}
 On a une solution non trivial $( B_\mathrm{I}, B_\mathrm{II}) \neq (0,0)$ si
 $$
k_\mathrm{II} = -  k_\mathrm{I} \text{cotan}(k_\mathrm{I}L).
 $$
On obtient la même condition que celle pour les modes anti-symétriques du puits symétrique fini ! Cela est logique, puisque les modes symétriques sont supprimés par la condition aux limites en $x=0$ (car $\vphi(x=0)=0$).
}

\item En déduire qu'il existe un nombre fini d'états liés que l'on caractérisera. Que se passe-t-il pour $V_{0}\to\infty$ ? Comparer au cas du puits infini et du puits fini symétrique.
\sol{
Posons $X=k_\mathrm{I}L$ et $Y= k_\mathrm{II}L$. Alors $Y=-X \cotan X$ d'après les conditions aux limites. L'après la définition de  $k_\mathrm{I}$ et $k_\mathrm{II}$, on a aussi
$$
X^2 + Y^2 = \frac{2m(V_0+E)}{\hbar^2} - \frac{2mE}{\hbar^2} =  \frac{2mV_0}{\hbar^2},
$$
qui est une constante (indépendant de l'énergie). Lorsque $V_0 \rightarrow \infty$, le nombre d’états liés tend vers l’infini, car le rayon du cercle défini par la condition ci-dessus tend vers l’infini.
}
\end{quest}
\end{exo}
%%%%%%%%%%%%%%%%%%%%%%%%%%%%%%%%%%%%%%%%%%%%%%%%%%%%%%%%%%%%%%%%%%%%%%%%%%%%%


%%%%%%%%%%%%%%%%%%%%%%%%%%%%%%%%%%%%%%%%%%%%%%%%%%%%%%%%%%%%%%%%%%%%%%%%%%%%%
\begin{exo}[PC][2]{Effet tunnel}
Une particule se déplace sur l'axe $(Ox)$ et est soumise à une barrière de potentiel de la forme
\begin{equation*}
V(x)=\begin{cases}
0 \mbox{ pour } x<0 \mbox{ (région I)},\\
V_{0} \mbox{ pour } 0<x<a \mbox{ (région II)},\\
0 \mbox{ pour } x>a \mbox{ (région III)}.
\end{cases}
\end{equation*}
On suppose que la particule arrive depuis la région I sur la barrière avec une énergie $0<E<V_0$.
\begin{quest}
\item Donner l'expression de la partie spatiale de la fonction d'onde $\vphi(x)$ dans chaque région en fonction de 6 constantes d'intégration.
\item  Justifier qu'il ne reste que 5 coefficients d'intégration et non 6. Préciser les relations de continuité entre les régions.
\item %On rappelle l'expression du courant de probabilité $J_{x}=\frac{i\hbar}{2m}\left(\psi\frac{\partial\psi^{\star}}{\partial x}-\psi^{\star}\frac{\partial\psi}{\partial x}\right)$ où $\psi^{\star}$ désigne le complexe conjugué de $\psi$.
Définir le coefficient de transmission $T$ et donner son expression.
\item La résolution du système précédent (on ne demande pas de le faire) conduit à
\begin{equation*}
T=\frac{1}{1+\frac{V_0^2}{4E(V_{0}-E)}\,\mathrm{sh}^{2}(k_{\rm{II}}a)}.
\end{equation*}
On donne $E=\SI{1}{eV}$, $h=\SI{6,62e-34}{J.s}$, $m=\SI{9,1e-31}{kg}$ et $V_{0}=\SI{2}{eV}$. Calculer le coefficient de transmission pour $a=\SI{0,5}{nm}$. On donnera tout d'abord en la justifiant une expression simplifiée du coefficient $T$.
\end{quest}
\end{exo}
%%%%%%%%%%%%%%%%%%%%%%%%%%%%%%%%%%%%%%%%%%%%%%%%%%%%%%%%%%%%%%%%%%%%%%%%%%%%%

%%%%%%%%%%%%%%%%%%%%%%%%%%%%%%%%%%%%%%%%%%%%%%%%%%%%%%%%%%%%%%%%%%%%%%%%%%%%%
\begin{exo}[PC][3][TD]{Microscopie à effet tunnel}
  \Notions{\item Équation de Schrödinger \task Onde stationnaire \task Effet tunnel}
  G. \textsc{Binnig} et H. \textsc{Rohrer} développèrent le microscope à effet tunnel chez IBM entre 1981 et 1985, ce qui leur valurent d'obtenir le prix Nobel de physique en 1986. La microscopie à effet tunnel consiste à relever le courant électrique $I$ circulant entre une pointe métallique et l'échantillon conducteur soumis à une tension $U$. Entre la pointe et l'échantillon, les électrons de conduction de masse $m$ traversent la couche d'air d'épaisseur $a$ par effet tunnel (Fig.~\ref{fig:tunnel}).

  \begin{fig}[tunnel]
    \centering
    \includegraphics[width=0.4\textwidth]{Microscope_effet_tunnel.svg.png}
    \capt{Schéma de principe du microscope à effet tunnel.}
  \end{fig}
En l'absence de tension, et en supposant que les électrons se comportent comme des particules libres dans l'échantillon, on modélise le potentiel ressenti par les électrons par une barrière symétrique :
\begin{equation*}
V(x)=\begin{cases}
0 \mbox{ pour } x<0 \mbox{ (région I)},\\
W \mbox{ pour } 0<x<a \mbox{ (région II)},\\
0 \mbox{ pour } x>a \mbox{ (région III)},
\end{cases}
\end{equation*}
où $W$ représente le travail d'extraction de l'échantillon conducteur.
La région I représente l'échantillon, la région II l'air, et la région III la pointe conductrice.
On suppose que l'électron arrive depuis la région I sur la barrière avec une énergie $0<E< W$. On rappelle que dans l'hypothèse d'une barrière épaisse $a\gg \delta$, le coefficient de transmission en probabilité $T(E)$ entre les régions I et III vérifie
% \begin{equation*}
% T(E) \simeq \frac{16E(W-E)}{E^2} e^{-\frac{2a}{\delta(E)}}, \qquad \text{avec } \delta(E) = \frac{\hbar}{\sqrt{2m(W-E)}}.
% \end{equation*}
\begin{equation*}
\ln T(E) \simeq -\dfrac{2a}{\delta(E)}, \qquad \text{avec } \delta(E) = \frac{\hbar}{\sqrt{2m(W-E)}}.
\end{equation*}
\begin{quest}
\item 
On donne  $h=\SI{6,6e-34}{J.s}$, l'énergie de l'électron $E=\SI{2}{eV}$, la masse de l'électron $m=\SI{9,1e-31}{kg}$, et $W=\SI{4.5}{eV}$. Calculer l'épaisseur de peau $\delta$ et le coefficient de transmission $T$ pour $a=\SI{0,5}{nm}$. 
\sol{
On calcule d'abord l'épaisseur de peau
$$
\delta = \frac{\hbar}{\sqrt{2m(W-E)}} = \SI{0.32}{nm}.
$$
On en déduit le coefficient de transmission
$$
\ln T(E) \simeq -\dfrac{2a}{\delta(E)} \Rar T(E) \simeq \e^{-2a/\delta} = \e^{-2\times 0.5/0.32} \simeq 0.04.
$$
}
\end{quest}
% On applique maintenant une tension $U$ entre la pointe et l'échantillon, de sorte que l'énergie des électrons vaut $E$ dans la région I et $E+eU$ dans la région III, où $e=\SI{1,6e-19}{C}$ est la charge élémentaire.
\begin{quest}
\item Justifier que la densité de courant électrique associé à \emph{un seul} électron de conduction sortant par la région III, d'énergie $E$, de vecteur d'onde $k_I$  et de densité de probabilité $\rho_I$ dans la région I, s'écrit  
\begin{equation*}
j_+(E) = -\frac{e\hbar k_I}{2m} T(E) \rho_I,
\end{equation*}
où $e=\SI{1,6e-19}{C}$ est la charge élémentaire.
\sol{
  La densité de courant de probabilité associée à une fonction d'onde $\psi$ est donnée par
  $$
  j = v_g \rho = \frac{\hbar k}{m} \rho,
  $$
  où $v_g = \frac{\hbar k}{m}$ est la vitesse de groupe de l'électron, et $\rho$ la densité de probabilité. La densité de courant de probabilité transmise, vaut donc par définition 
  $$
  j_t = T(E) j_i = T(E) \frac{\hbar k}{m} \rho_I.
  $$
  La densité de courant électrique est reliée à la densité de courant de probabilité par la relation $j_+ = q j_t$, où $q$ désigne la charge de la particule. Ici, $q=e$. Donc
 \eqbox{
    j_+(E) = \frac{e\hbar k}{m} T(E) \rho_I.
 }
}
\item Résoudre l'équation de Schrödinger indépendante du temps dans la région I pour déterminer la fonction d'onde spatiale $\vphi_I(x)$ et les valeurs possibles de $k_I$. On supposera que l'échantillon est de longueur~$L$ et on imposera des conditions aux limites périodiques. En déduire la densité de probabilité~$\rho_I$.
% \sol{
% L'équation de Schrödinger indépendante du temps dans la région I s'écrit
% $$
% \frac{\d^2 \vphi_I}{\d x^2} + \frac{2mE}{\hbar^2} \vphi_I = 0.  
% $$
% La solution générale est
% $$
% \vphi_I(x) = A \cos(k_I x) + B \sin(k_I x), \quad \text{avec } k_I = \sqrt{\frac{2mE}{\hbar^2}}.
% $$
% Comme l'échantillon est de longueur finie $L$, et que l'énergie d'extraction $W$ est très grande devant l'énergie de l'électron $E$, tout se passe comme si l'électron était confiné dans un puits de potentiel infini. On impose les condition aux limites rigides de continuité de la fonction d'onde 
% $$\vphi_I(-L) = \vphi_I(0) = 0,$$
% ce qui conduit à $A=0$ et $B$ quelconque avec 
% \eqbox{
% k_I = n \frac{\pi}{L}, \quad n \in \N^\star.
% }
% avec $n \in \N^\star$. La fonction d'onde s'écrit donc
% $$
% \vphi_I(x) = B \sin\left( n \frac{\pi}{L} x \right).
% $$
% La densité de probabilité s'obtient par 
% $$\rho_I = |\vphi_I(x)|^2 = |B|^2 \sin^2\left( n \frac{\pi}{L} x \right).$$ La constante $B$ s'obtient par la normalisation de la fonction d'onde, soit
% $$
% \int_0^L \rho_I \d x = 1 \Rar |B|^2 \int_0^L \sin^2\left( n \frac{\pi}{L} x \right) \d x = 1 \Rar |B|^2 = \frac{2}{L}.
% $$
% On en déduit
% \eqbox{
%   \rho_I(x) = \frac{2}{L} \sin^2\left( n \frac{\pi}{L} x \right). 
% }
% }
\sol{
L'équation de Schrödinger indépendante du temps dans la région I s'écrit
$$
\frac{\d^2 \vphi_I}{\d x^2} + \frac{2mE}{\hbar^2} \vphi_I = 0.  
$$
La solution générale est
$$
\vphi_I(x) = A \e^{ i k_I x} + B \e^{- i k_I x}, \quad \text{avec } k_I = \sqrt{\frac{2mE}{\hbar^2}}.
$$
Comme seuls les électrons qui se  déplacent vers les $x$ croissants contribuent au courant vers la région  III, on impose $B=0$, soit $\vphi_I(x) = A \e^{ i k_I x}$.
Les conditions aux limites périodiques imposent  alors
$$\vphi_I(-L) = \vphi_I(0),$$
ce qui conduit 
\eqbox{
k_I = n \frac{2\pi}{L},
}
avec $n \in \N^\star$. La fonction d'onde s'écrit donc
$$
\vphi_I(x) = A \e^{ i n \frac{2\pi}{L} x}.
$$
La densité de probabilité s'obtient par 
$$\rho_I = |\vphi_I(x)|^2 = |A|^2.$$ La constante $A$ s'obtient par la normalisation de la fonction d'onde, soit
$$
\int_0^L \rho_I \d x = 1 \Rar |A|^2 \int_0^L \d x = 1 \Rar |A|^2 = \frac{1}{L}.
$$
On en déduit
\eqbox{
  \rho_I(x) = \frac{1}{L}. 
}
}
\end{quest}
En réalité, un métal contient un grand nombre d'électrons de conduction, dont les énergies sont réparties selon la statistique de Fermi-Dirac. Dans la limite des basses températures, tous les états quantiques de l'échantillon sont occupés par un électron sous une énergie $\varepsilon_F$, appelée énergie de Fermi, tandis que les états au-dessus d'énergie supérieure à $\varepsilon_F$ sont vides. Dans la région III, la tension $U$ appliquée entre la pointe et l'échantillon modifie l'énergie de Fermi de l'échantillon en $\veps_F + eU$.
\begin{quest}
\item Déterminer la densité de courant $\d J_+$ issue des électrons de vecteur d'onde compris entre $k$ et $k+\d k$, puis montrer que le courant électrique  de la région I vers la région III s'écrit
\begin{equation*}
J_+(E) = -\frac{2e}{h} \int_0^{\varepsilon_F} T(E)\d E,
\end{equation*}
\sol{
Chaque état quantique peut être occupé par deux électrons de spins opposés, il y a donc deux états quanstiques par intervalle $\Delta k = \frac{2\pi}{L}$. Le nombre d'états quantiques entre $k_I$ et $k_I + \d k$ vaut donc
$$
\d N  = 2 \frac{\d k}{\Delta k} = \frac{L}{\pi} \d k.
$$
Le courant électrique élémentaire associé à ces états quantiques est
$$
\d J_+ = j_+(E) \d N = -\frac{e\hbar k}{m} T(E) \rho_I \frac{L}{\pi} \d k.
 = -\frac{e\hbar k}{m \pi} T(E) \d k.
$$
En intégrant sur tous les états occupés, soit de l'énergie $0$ à l'énergie de Fermi $\veps_F$, on obtient
$$
J_+ = \int_0^{k_F} \d J_+ = -\frac{e}{\hbar \pi} \int_0^{k_F} T(E) \frac{\hbar^2 k}{m}\d k,
$$
soit 
\eqbox{
 J_+(E) = -\frac{2e}{h} \int_0^{\veps_F} T(E) \d E.
}
}
\item En considérant les électrons se déplaçant de la région III vers la région I, déterminer l'intensité du courant électrique total $I$, et en déduire l'expression de Landauer-Büttiker de la conductance
\begin{equation*}
G = G_0 T(\veps_F), \quad \text{avec } G_0 = \frac{2e^2}{h}
\end{equation*}
 le quantum de conductance. Évaluer numériquement $G_0$ et $G$ pour une énergie de Fermi $\veps_F = \SI{2}{eV}$.
 \sol{
  Comme les porteurs de charge (électrons) sont négatifs, le courant électrique total vers les $x$ croissants s'écrit
  $$
I = J_-(E) - J_+(E) = \frac{2e}{h} \int_{0}^{\veps_F + eU} T(E) \d E - \frac{2e}{h} \int_{0}^{\veps_F} T(E) \d E = \frac{2e}{h} \int_{\veps_F}^{\veps_F + eU} T(E) \d E.
  $$
  Dans la limite des faibles tensions $eU \ll \veps_F$, on peut considérer que $T(E) \simeq T(\veps_F)$ dans l'intervalle d'intégration, et donc
  \eqbox{
    I \simeq \frac{2e^2}{h} T(\veps_F) U.
  }
  La loi d'Ohm microscopique $I = G U$ permet d'identifier la conductance
  \eqbox{
    G = G_0 T(\veps_F), \quad \text{avec } G_0 = \frac{2e^2}{h}.
  }
  Numériquement, on trouve
  $$
  G_0 \simeq \SI{7.75e-5}{S}, \quad G \simeq \SI{7.75e-5}{S}.
  $$
 }
\item Un microscope à effet tunnel permet de scanner la surface d'un matériau pour en obtenir un profil topographique $z_0(x)$ de l'échantillon à l'échelle atomique. Il existe deux méthodes de fonctionnement :  à altitude $z$ au-dessus de l'échantillon imposée, ou à courant tunnel $I$ imposé. Expliquer le principe de ces deux méthodes.
\sol{
  À énergie de Fermi $\veps_F$ et tension $U$ donnés, l'intensité du courant électrique total s'écrit 
  $$
  I(a) \simeq \frac{2e^2}{h} V e^{-2a/\delta}.
  $$
  L'intensité du courant $I$ est donc constante si la hauteur entre la pointe et l'échantillon $a$ est constante.
  \begin{enumerate}
    \item À altitude $z$ imposée : la distance $a$ entre la pointe et l'échantillon  vaut alors $a=z-z_0(x)$. On mesure le courant tunnel 
    $$
    I(x) = I(z - z_0(x)) \simeq \frac{2e^2}{h} V e^{-2(z-z_0(x))/\delta} \propto \e^{2z_0(x)/\delta},
    $$
    pour en déduire le profil topographique $z_0(x)$ de l'échantillon.
    \item À courant tunnel $I$ imposé : on ajuste la hauteur $z$ de la pointe au-dessus de l'échantillon grâce à un système de rétroaction piezoélectrique pour que la distance $a$ soit constante. On en déduit le profil topographique $z_0(x)$ de l'échantillon à partir de la mesure de la hauteur $z(x) = a + z_0(x)$.
  \end{enumerate}
}
%Soit $n$ la densité particulaire en électron dans l'échantillon conducteur, $e$ la charge élémentaire, et $\tau$
% On donnera tout d'abord en la justifiant une expression simplifiée du coefficient $T$.
\end{quest}
\end{exo}
%%%%%%%%%%%%%%%%%%%%%%%%%%%%%%%%%%%%%%%%%%%%%%%%%%%%%%%%%%%%%%%%%%%%%%%%%%%%%
